\begin{center}
{\Large\bfseries\color{ink} Differenciálegyenletek}\\[5pt]
{\large\bfseries 1. előadás}\\[4pt]
{\small 2026. ősz}\qquad\textcolor{graytext}{\small Az eredeti kézirat dátuma: 2019. 09. 11. (1)}
\end{center}

\begin{matbox}
\textbf{Előadó:} Szili László \hfill \textbf{Időpont:} Kedd 8:15--9:45

\textbf{Honlapon:}
\begin{itemize}[leftmargin=1.4em]
\item \textbf{Irodalom:} Simon L. Péter, Tóth János: d.e. (2008); Pontrjagin: közönséges differenciálegyenletek
\item \textbf{Követelmény:} vizsgajegy: az előadáson elhangzottak.
\item \textbf{Szimbolikus programcsomagok:} Mathematica (Wolfram), Maple (Differenciálegyenlet lehetőségek), Derive.
\end{itemize}
\end{matbox}

\section{Bevezetés: Algebrai és differenciálegyenletek}

\fej{Emlékeztető: Algebrai egyenletek}
Adott egy $f\colon\mathbb R\to\mathbb R$ függvény; keressük az
$x\in D_f$ számot, amelyre
\[
 f(x)=0.
\]

\fej{Differenciálegyenlet}
Az is egy feladat, de az \emph{ismeretlen most függvény}.

\medskip
\begin{center}
\begin{tikzpicture}[
  base/.style={draw=ink!45,rounded corners=2pt,align=center,inner sep=5pt},
  node distance=13mm]
\node[base,draw=accent,thick,fill=soft] (top) {\textbf{Differenciálegyenlet (DE)}};
\node[base,below left=12mm and 4mm of top] (ode)
  {\textbf{Közönséges DE}\\[-1pt]\footnotesize ismeretlen: $\mathbb R\to\mathbb R$ függvény};
\node[base,below right=12mm and 4mm of top] (pde)
  {\textbf{Parciális DE}\\[-1pt]\footnotesize ismeretlen: többváltozós függvény};
\draw[ink,semithick] (top.south) -- ++(0,-5mm) -| (ode.north);
\draw[ink,semithick] (top.south) -- ++(0,-5mm) -| (pde.north);
\end{tikzpicture}
\end{center}

\fej{Közönséges differenciálegyenletek fajtái}
\begin{enumerate}[label=\arabic*)]
\item \textbf{Explicit elsőrendű:}
\[
 x'(t)=f\bigl(t,x(t)\bigr).
\]
\item \textbf{Implicit elsőrendű:} a derivált nincs kifejezve:
\[
 F\bigl(t,x(t),x'(t)\bigr)=0.
\]
\item \textbf{Magasabbrendű differenciálegyenletek:} magasabb rendű n>2 deriváltak.
\item \textbf{Differenciálegyenlet-rendszer.}
\end{enumerate}

\fej{Parciális differenciálegyenlet -- példa}
\[
 \frac{\partial^2 u(x,y)}{\partial x^2}
 +\frac{\partial^2 u(x,y)}{\partial y^2}=0.
\]
Ez a \emph{Laplace-egyenlet} (ld.: Analízis alkalmazásai 2022-23 tanév, ősz, 5. előadás).


\section{Bevezető példák}

\subsection*{1. példa: adott függvény primitív függvényei}
\begin{matbox}
\[
 \boxed{x'(t)=h(t)},\qquad t\in I\subseteq\mathbb R,
 \quad I\text{ nyílt intervallum}.
\]
Adott $h\colon I\to\mathbb R$, ahol $h\in C(I)$.
\end{matbox}

\textbf{Megoldás:} az $x(t)$ megoldások éppen $h$ primitív függvényei.

\subsection*{2. példa: $x'=x$ és a kezdetiérték-probléma}

Határozzuk meg az összes olyan $x\colon I\to\mathbb R$ függvényt, amelyre
\begin{matbox}
\[
 \boxed{x'(t)=x(t)}\qquad (t\in I,\ I\text{ nyílt intervallum}).
\]
\end{matbox}

\textbf{Megoldások:} $I\subseteq\mathbb R$ tetszőleges nyílt intervallumon
\[
 \boxed{x(t)=c\,e^t}\qquad (t\in I),\quad \forall c\in\mathbb R.
\]

\textbf{Van más megoldás?} Nincs!

\textbf{Ötlet:} Tegyük fel, hogy $\varphi\colon I\to\mathbb R$
megoldás, tehát $\varphi'(t)=\varphi(t)$, és tekintsük a
$\dfrac{\varphi(t)}{e^t}$ függvényt (Lagrange-középértéktétel):
\begin{align*}
\left(\frac{\varphi(t)}{e^t}\right)'
 &=\frac{\varphi'(t)e^t-\varphi(t)e^t}{e^{2t}}\\
 &=0
 \quad\Longrightarrow\quad
 \exists c\in\mathbb R:\quad
 \frac{\varphi(t)}{e^t}=c
 \quad(\forall t\in I).
\end{align*}
Következésképpen $\varphi(t)=ce^t$.

\newpage
\fej{Szemléltetés: megoldásgörbék}
\begin{center}
\begin{tikzpicture}[scale=0.94]
\begin{scope}
\clip (-2.8,-2.0) rectangle (3.0,2.9);
\foreach \c in {0.15,0.3,0.62,1.0,-0.13,-0.27,-0.52,-0.90}{
\draw[curve] plot[domain=-2.8:2.9,samples=70] (\x,{\c*exp(\x)});
}
\draw[guide] (0,1.13)--(0.60,1.13)--(0.60,0);
\fill[ink] (0.60,1.13) circle(1.8pt);
\end{scope}
\draw[axis] (-2.8,0)--(3,0) node[below] {$t$};
\draw[axis] (0,-2.0)--(0,2.9) node[above] {$x$};
\node[left] at (0,1.13) {$\xi$};
\node[below] at (0.60,0) {$\tau$};
\node at (2.18,2.50) {$c>0$};
\node at (2.25,-1.58) {$c<0$};
\end{tikzpicture}

Kezdetiérték-probléma (k.é.p.)
\[
\boxed{\begin{cases}
 x'(t)=x(t),\\
 x(\tau)=\xi.
\end{cases}}
\]
\end{center}

\subsection*{3. példa: $x'=\sqrt{|x|}$}

Az összes megoldást keressük:
\begin{matbox}
\[
\boxed{x'(t)=\sqrt{\lvert x(t)\rvert}},\qquad t\in\mathbb R.
\]
\end{matbox}

\textbf{Megoldásvázlat:} a kezdetiérték-probléma
\[
\begin{cases}
 x'(t)=\sqrt{|x(t)|},\\
 x(\tau)=\xi.
\end{cases}
\]
Három esetet különböztetünk meg.

\fej{(i) $\xi>0$, a vizsgált szakaszon $x(t)>0$}
Tegyük fel, hogy $\varphi(t)>0$ $\forall$ minden $t\in I$ esetén. Ekkor
\[
\varphi'(t)=\sqrt{\varphi(t)},
\]
így
\[
\frac{\varphi'(t)}{\sqrt{\varphi(t)}}
 =\left(2\sqrt{\varphi(t)}\right)'=1.
\]

\begin{center}
\begin{tikzpicture}[scale=0.67]
\draw[axis] (-2.2,0)--(2.6,0) node[below] {$t$};
\draw[axis] (0,-0.75)--(0,2.35) node[above] {$x$};
\draw[curve] plot[domain=-1.65:1.36,samples=45] (\x,{.25*(\x+1.65)^2});
\draw[curve] plot[domain=-0.7:2.2,samples=45] (\x,{.25*(\x+0.70)^2});
\draw[guide] (0,1)--(0.35,1)--(0.35,0);
\fill[ink] (0.35,1) circle(1.8pt);
\node[left] at (0,1) {$\xi$};
\node[below] at (0.35,0) {$\tau$};
\end{tikzpicture}
\figlabel{Pozitív kezdeti érték}
\end{center}

\fej{(ii) $\xi<0$, a vizsgált szakaszon $x(t)<0$}
A jegyzet a negatív tartományban futó, a $t$-tengelyt elérő
megoldásgörbéket szemlélteti.

\begin{center}
\begin{tikzpicture}[scale=0.67]
\draw[axis] (-2.2,0)--(2.6,0) node[below] {$t$};
\draw[axis] (0,-2.30)--(0,0.9) node[above] {$x$};
\foreach \a in {0.45,1.35,2.35}{
\draw[curve] plot[domain=-2.1:\a,samples=45] (\x,{-.25*(\a-\x)^2});
\draw[curve] (\a,0)--(2.3,0);
}
\draw[guide] (0,-0.46)--(1,-0.46)--(1,0);
\fill[ink] (1,-0.46) circle(1.8pt);
\node[left] at (0,-0.46) {$\xi$};
\node[above] at (1,0.06) {$\tau$};
\end{tikzpicture}
\figlabel{Negatív kezdeti érték}
\end{center}

\fej{(iii) $\xi=0$, $x(\tau)=0$}
\[
\varphi(t)=0\qquad\text{megoldás.}
\]
A kézirat ezenkívül több, a $t$-tengelyhez csatlakozó
megoldásgörbét is ábrázol.

\begin{center}
\begin{tikzpicture}[scale=0.67]
\draw[axis] (-2.2,0)--(2.6,0) node[below] {$t$};
\draw[axis] (0,-1.65)--(0,2.05) node[above] {$x$};
\draw[curve] (-2.1,0)--(2.3,0);
\draw[curve] plot[domain=-2.1:-0.25,samples=40] (\x,{-.40*(-0.25-\x)^2});
\draw[curve] (-0.25,0)--(0.65,0);
\draw[curve] plot[domain=0.65:2.4,samples=40] (\x,{.50*(\x-0.65)^2});
\draw[curve] plot[domain=-1.75:0.55,samples=40] (\x,{-.26*(0.55-\x)^2});
\draw[curve] (0.55,0)--(1.15,0);
\draw[curve] plot[domain=1.15:2.4,samples=40] (\x,{.45*(\x-1.15)^2});
\fill[ink] (0.55,0) circle(1.8pt);
\node[below] at (0.55,0) {$\tau$};
\node[left] at (0,0) {$\xi=0$};
\end{tikzpicture}
\figlabel{Zérus kezdeti érték}
\end{center}

\fej{Szemléltetés: „parabolaágak”}
\[
\boxed{\displaystyle
\varphi(t)=\left(\frac{t+c}{2}\right)^2},\qquad t>-c.
\]
\begin{center}
\begin{tikzpicture}[scale=0.91]
\begin{scope}
\clip (-2.8,-0.2) rectangle (3.5,2.5);
\foreach \a in {-2.45,-0.95,0.65}{
\draw[curve] plot[domain=\a:3.3,samples=70] (\x,{.25*(\x-\a)^2});
}
\end{scope}
\draw[axis] (-2.8,0)--(3.5,0) node[below] {$t$};
\draw[axis] (0,-0.15)--(0,2.5) node[above] {$x$};
\draw[guide] (0,1)--(2.65,1)--(2.65,0);
\fill[ink] (2.65,1) circle(1.8pt);
\node[left] at (0,1) {$\xi$};
\node[below] at (2.65,0) {$\tau$};
\node[above right] at (1.50,2.5) {„parabolaágak”};
\end{tikzpicture}
\end{center}
$\tau$ meghatározható


\subsection*{4. példa: a szignumfüggvény}
\begin{matbox}
\[
\boxed{x'(t)=\operatorname{sgn}(t)},\qquad t\in(-1,1).
\]
\end{matbox}

\begin{center}
\begin{tikzpicture}[scale=0.92]
\draw[axis] (-2,0)--(2.25,0) node[below] {$t$};
\draw[axis] (0,-1.45)--(0,1.45) node[above] {$x$};
\draw[curve,very thick] (-1.7,-1)--(0,-1);
\draw[curve,very thick] (0,1)--(1.7,1);
\draw[fill=white,draw=accent,semithick] (0,-1) circle(2.5pt);
\draw[fill=white,draw=accent,semithick] (0,1) circle(2.5pt);
\fill[ink] (0,0) circle(2.5pt);
\node[left] at (0,1) {$1$};
\node[right] at (0,-1) {$-1$};
\end{tikzpicture}
\figlabel{A $\operatorname{sgn}(t)$ függvény grafikonja.}
\end{center}

Ennek \textbf{nincs megoldása}, hiszen a szignumfüggvénynek
$(-1,1)$-en nincs primitív függvénye
\emph{(nem Darboux-tulajdonságú)}.

\begin{matbox}
\textbf{Megjegyzés:} A deriváltfüggvény mindig Darboux-tulajdonságú.

Primitív függvény létezésének \textbf{szükséges feltétele}, hogy
az adott függvény Darboux-tulajdonságú legyen.
\end{matbox}

